% GENERATED FILE. Edit canonical lessons, then run npm run generate:books.
% canonical-source-hash: fnv1a64:286f5df56f1c012b
% canonical-lessons: RU-C37-neighbour, RU-C37-guest, RU-C37-teacher, RU-C37-pupil, RU-C37-work

\chapter{Neighbour, Guest, Teacher, Pupil, Work}
\label{ch:ru-neighbour-guest-teacher-pupil-work}
\hlchaptermodality{37}

\begin{chapteropening}
\textbf{By the end of this chapter:} \emph{I can name a neighbour, a guest, a teacher, a pupil and work.}

Complete the last lesson of chapter 37: I can name a neighbour, a guest, a teacher, a pupil and work.
\end{chapteropening}

\section[\texorpdfstring{\ru{сосед} (soséd)}{сосед (soséd)}]{\ru{сосед} (soséd) — a neighbour}

\label{lesson:RU-C37-neighbour}



\begin{quote}
Before the new one: say the Russian for a grandson, then the Russian for a granddaughter.
\end{quote}

\subsection*{You'll want to know: \ru{сосед}}
\textbf{\ru{сосед}} (\emph{soséd}) — \textquotedblleft{}a neighbour\textquotedblright{}.

\begin{cousinweb}
\textbf{\ru{сосед}} $\leftarrow$ Proto-Slavic \emph{*sǫsědŭ} — inherited: \emph{*sǫ-}, \textquotedblleft{}together\textquotedblright{}, with the root of \emph{\ru{сидеть}}, \textquotedblleft{}to sit\textquotedblright{}, so a neighbour is \textquotedblleft{}one who sits beside you\textquotedblright{}. Vasmer compares Sanskrit \emph{saṃsád}, \textquotedblleft{}assembly\textquotedblright{}.
\end{cousinweb}

\begin{grammarlens}[title={What you've built}]
One more word for this chapter.
\end{grammarlens}

\subsection*{Your turn}
\begin{itemize}
\raggedright
  \item \emph{Say it:} \emph{soséd}
  \item \emph{Say it:} \emph{soséd}, once more
  \item \emph{Say it:} say \emph{soséd}
  \item \emph{Recall:} say the Russian for a grandson, then the Russian for a granddaughter, then say \emph{soséd} again
\end{itemize}

\begin{tcolorbox}[breakable,colback=teal!4,colframe=teal!35!black,title={Before you move on}]
What does \ru{сосед} mean? (\textquotedblleft{}A neighbour\textquotedblright{}.) Say it once more.
\end{tcolorbox}
\section[\texorpdfstring{\ru{гость} (gost)}{гость (gost)}]{\ru{гость} (gost) — a guest}

\label{lesson:RU-C37-guest}



\begin{quote}
Before the new one: say the Russian for a granddaughter, then the Russian for a neighbour.
\end{quote}

\subsection*{You'll want to know: \ru{гость}}
\textbf{\ru{гость}} (\emph{gost}) — \textquotedblleft{}a guest\textquotedblright{}.

\begin{cousinweb}
\textbf{\ru{гость}} $\leftarrow$ Proto-Slavic \emph{*gostĭ} — inherited. Vasmer relates it to Gothic \emph{gasts} and Latin \emph{hostis}, \textquotedblleft{}stranger, enemy\textquotedblright{}, and rejects a loan from Germanic: English \emph{guest} is a cousin, from PIE \emph{*g\textsuperscript{h}óstis}. In Old Russian a \emph{\ru{гость}} was also a travelling merchant (Vasmer).
\end{cousinweb}

\begin{grammarlens}[title={What you've built}]
One more word for this chapter.
\end{grammarlens}

\subsection*{Your turn}
\begin{itemize}
\raggedright
  \item \emph{Say it:} \emph{gost}
  \item \emph{Say it:} \emph{gost}, once more
  \item \emph{Say it:} say \emph{gost}
  \item \emph{Recall:} say the Russian for a granddaughter, then the Russian for a neighbour, then say \emph{gost} again
\end{itemize}

\begin{tcolorbox}[breakable,colback=teal!4,colframe=teal!35!black,title={Before you move on}]
What does \ru{гость} mean? (\textquotedblleft{}A guest\textquotedblright{}.) Say it once more.
\end{tcolorbox}
\section[\texorpdfstring{\ru{учитель} (uchítel)}{учитель (uchítel)}]{\ru{учитель} (uchítel) — a teacher}

\label{lesson:RU-C37-teacher}



\begin{quote}
Before the new one: say the Russian for a neighbour, then the Russian for a guest.
\end{quote}

\subsection*{You'll want to know: \ru{учитель}}
\textbf{\ru{учитель}} (\emph{uchítel}) — \textquotedblleft{}a teacher\textquotedblright{}.

\begin{cousinweb}
\textbf{\ru{учитель}} is \emph{\ru{учить}}, \textquotedblleft{}to teach\textquotedblright{}, with \emph{-\ru{тель}}, \textquotedblleft{}one who does it\textquotedblright{}; Old Church Slavonic already had it. The verb goes back to Proto-Slavic \emph{*učiti}, \textquotedblleft{}to accustom\textquotedblright{} (Vasmer).
\end{cousinweb}

\begin{grammarlens}[title={What you've built}]
One more word for this chapter.
\end{grammarlens}

\subsection*{Your turn}
\begin{itemize}
\raggedright
  \item \emph{Say it:} \emph{uchítel}
  \item \emph{Say it:} \emph{uchítel}, once more
  \item \emph{Say it:} say \emph{uchítel}
  \item \emph{Recall:} say the Russian for a neighbour, then the Russian for a guest, then say \emph{uchítel} again
\end{itemize}

\begin{tcolorbox}[breakable,colback=teal!4,colframe=teal!35!black,title={Before you move on}]
What does \ru{учитель} mean? (\textquotedblleft{}A teacher\textquotedblright{}.) Say it once more.
\end{tcolorbox}
\section[\texorpdfstring{\ru{ученик} (ucheník)}{ученик (ucheník)}]{\ru{ученик} (ucheník) — a pupil}

\label{lesson:RU-C37-pupil}



\begin{quote}
Before the new one: say the Russian for a guest, then the Russian for a teacher.
\end{quote}

\subsection*{You'll want to know: \ru{ученик}}
\textbf{\ru{ученик}} (\emph{ucheník}) — \textquotedblleft{}a pupil\textquotedblright{}.

\begin{cousinweb}
\textbf{\ru{ученик}} is \emph{\ru{учить}}, \textquotedblleft{}to teach, to learn\textquotedblright{}, with \emph{-\ru{ник}}, a suffix for a person: the one who learns. The verb goes back to Proto-Slavic \emph{*učiti}, \textquotedblleft{}to accustom\textquotedblright{} (Vasmer).
\end{cousinweb}

\begin{grammarlens}[title={What you've built}]
One more word for this chapter.
\end{grammarlens}

\subsection*{Your turn}
\begin{itemize}
\raggedright
  \item \emph{Say it:} \emph{ucheník}
  \item \emph{Say it:} \emph{ucheník}, once more
  \item \emph{Say it:} say \emph{ucheník}
  \item \emph{Recall:} say the Russian for a guest, then the Russian for a teacher, then say \emph{ucheník} again
\end{itemize}

\begin{tcolorbox}[breakable,colback=teal!4,colframe=teal!35!black,title={Before you move on}]
What does \ru{ученик} mean? (\textquotedblleft{}A pupil\textquotedblright{}.) Say it once more.
\end{tcolorbox}
\section[\texorpdfstring{\ru{работа} (rabóta)}{работа (rabóta)}]{\ru{работа} (rabóta) — work, a job}

\label{lesson:RU-C37-work}



\begin{quote}
Before the new one: say the Russian for a teacher, then the Russian for a pupil.
\end{quote}

\subsection*{You'll want to know: \ru{работа}}
\textbf{\ru{работа}} (\emph{rabóta}) — \textquotedblleft{}work, a job\textquotedblright{}.

\begin{cousinweb}
\textbf{\ru{работа}} $\leftarrow$ Proto-Slavic \emph{*orbota}: Old Russian \emph{\ru{робота}} meant \textquotedblleft{}work, slavery\textquotedblright{} (Vasmer). It is built on the root of \emph{\ru{раб}}, \textquotedblleft{}slave\textquotedblright{}, and Vasmer relates it to Gothic \emph{arbaiþs}, \textquotedblleft{}hardship\textquotedblright{}, and German \emph{Arbeit}, \textquotedblleft{}work\textquotedblright{}. Its \emph{\ru{ра}-} is the Church Slavonic shape; East Slavic dialects say \emph{\ru{робота}}.
\end{cousinweb}

\begin{grammarlens}[title={What you've built}]
One more word for this chapter.
\end{grammarlens}

\subsection*{Your turn}
\begin{itemize}
\raggedright
  \item \emph{Say it:} \emph{rabóta}
  \item \emph{Say it:} \emph{rabóta}, once more
  \item \emph{Say it:} say \emph{rabóta}
  \item \emph{Recall:} say the Russian for a teacher, then the Russian for a pupil, then say \emph{rabóta} again
\end{itemize}

\begin{tcolorbox}[breakable,colback=teal!4,colframe=teal!35!black,title={Before you move on}]
What does \ru{работа} mean? (\textquotedblleft{}Work, a job\textquotedblright{}.) Say it once more.
\end{tcolorbox}
